\section{Reduction of the virtual root functions}
\subsection{Collecting the coefficients}
A compact expression need not eliminate the $G$ functions. It is useful
first to combine their coefficients, before expanding any function into
classical polylogarithms. For the six root words in the preceding table,
the shuffle product gives
\begin{equation}
 G_{0,r,-1}=\ln X\,G_{r,-1}-G_{r,0,-1}-G_{r,-1,0}.
 \label{eq:comp-root-compact-shuffle}
\end{equation}
Thus $G_{0,r,-1}$ is redundant. After this substitution, the coefficients
of $G_{r,0,-1}$ and $G_{r,-1,0}$ are proportional with ratio $1:-3$.
Define
\begin{align}
 p_F(X)&=7X^4+18X^3-144X^2-512X-256,\\
 p_K(X)&=3X^4+10X^3-48X^2-192X-96,\\
 p_W(X)&=3X^4+8X^3-60X^2-216X-112,\\
 p_Z(X)&=X^4+2X^3-24X^2-80X-48,\\
 p_0(X)&=X^4-2X^3-48X^2-128X-160.
 \label{eq:comp-root-compact-polynomials}
\end{align}
The contribution $C_{\mathrm{six}}$ of these six original words to
$\widehat C_{V+\mathrm{UV}}$ becomes
\begin{equation}
 \begin{aligned}
 C_{\mathrm{six}}(X)=\frac{Y}{X^4(X+4)}\bigl\{&
 3p_FG_{r,-1,-1}+p_K(G_{r,0,-1}-3G_{r,-1,0})\\
 &-6p_WG_{-4,r,-1}-8p_ZG_{r,-2,-1}
       +3p_0\ln X\,G_{r,-1}\bigr\}.
 \end{aligned}
 \label{eq:comp-root-compact-G}
\end{equation}
All coefficients are obtained by exact rational collection of the original
table. This removes one weight-three function and gives the other five a
common denominator. The $G$ notation keeps this relation shorter than its
fully expanded classical form. The identities below provide alternative
representations for evaluation and further algebraic reductions.
